\documentclass[a4paper,13pt,3p,twoside]{report}
\usepackage{scrextend}
\changefontsizes{13pt}
\usepackage[utf8]{vietnam}
\usepackage[top=2cm, bottom=2cm, left=3.5cm, right=2.5cm]{geometry}

\usepackage{graphicx} % Cho phép chèn hỉnh ảnh
\usepackage{fancybox} % Tạo khung box
\usepackage{indentfirst} % Thụt đầu dòng ở dòng đầu tiên trong đoạn
\usepackage{amsthm} % Cho phép thêm các môi trường định nghĩa
\usepackage{latexsym} % Các kí hiệu toán học
\usepackage{amsmath} % Hỗ trợ một số biểu thức toán học
\usepackage{amssymb} % Bổ sung thêm kí hiệu về toán học
\usepackage{amsbsy} % Hỗ trợ các kí hiệu in đậm
\usepackage{times} % Chọn font Time New Romans
\usepackage{array} % Tạo bảng array
\usepackage{enumitem} % Cho phép thay đổi kí hiệu của list
\usepackage{subfiles} % Chèn các file nhỏ, giúp chia các chapter ra nhiều file hơn
\usepackage{titlesec} % Giúp chỉnh sửa các tiêu đề, đề mục như chương, phần,..
\usepackage{titletoc}
\usepackage{chngcntr} % Dùng để thiết lập lại cách đánh số caption,..
\usepackage{pdflscape} % Đưa các bảng có kích thước đặt theo chiều ngang giấy
\usepackage{afterpage}
\usepackage[ruled,vlined]{algorithm2e}  % Hỗ trợ viết các giải thuật
\DontPrintSemicolon                      % bỏ dấu ';' cuối dòng cho gọn
\newcommand{\algcommentfont}[1]{\textcolor{black!55}{\itshape #1}} % chú thích xám nhạt, in nghiêng
\SetCommentSty{algcommentfont}           % chú thích bớt nổi để đỡ rối mắt
\SetAlgoNlRelativeSize{-1}               % số dòng nhỏ hơn một bậc
\setlength{\algomargin}{1.8em}           % thụt lề thân thuật toán rõ ràng
\SetArgSty{textnormal}                   % tham số không in nghiêng đậm, dễ đọc
\usepackage{capt-of} % Cho phép sử dụng caption lớn đối với landscape page
\usepackage{multirow} % Merge cells
\usepackage{booktabs} % Bảng chuyên nghiệp: \toprule \midrule \bottomrule
\renewcommand{\arraystretch}{1.3} % giãn chiều cao hàng, tránh chữ chạm viền bảng
\usepackage{fancyhdr} % Cho phép tùy biến header và footer
% \usepackage[natbib,backend=biber,style=ieee]{biblatex} % Giúp chèn tài liệu tham khảo
\usepackage{appendix}

\usepackage[font=small,labelfont=bf]{caption}

\usepackage{listings}
\usepackage{float}
\usepackage{subcaption}
\usepackage{xurl}
\usepackage{tikz} % Vẽ biểu đồ use case
\usetikzlibrary{positioning,shapes.geometric,arrows.meta,calc,fit,backgrounds,shapes.multipart}

% Khung giữ chỗ cho ảnh chụp màn hình giao diện. Sinh viên thay khối này bằng
% \includegraphics khi đã chụp được ảnh thực tế của ứng dụng đang chạy.
\newcommand{\uiscreenshot}[1]{%
  \fbox{\begin{minipage}[c][4.3cm][c]{0.82\textwidth}\centering\itshape
  [\,Ảnh chụp màn hình\,]\\[3pt] #1\\[4pt]
  \footnotesize\upshape(chèn tệp ảnh tương ứng trong thư mục \texttt{Hinhve/})
  \end{minipage}}}

\usepackage[nonumberlist, nopostdot, nogroupskip, acronym]{glossaries}
\usepackage{glossary-superragged}
\setglossarystyle{superraggedheaderborder}
\usepackage{setspace}
\usepackage{parskip}

% package content table
\usepackage{tocbasic}

\usepackage{blindtext}
\usepackage{outlines} % môi trường \begin{outline} dùng trong Phụ lục A (template thiếu gói này)


% ===================================================

\renewcommand{\bibname}{Danh_sach_tai_lieu_tham_khao} 
\usepackage[backend=bibtex,style=ieee]{biblatex}  %backend=biber is 'better'

\renewcommand\appendixname{PHỤ LỤC}
\renewcommand\appendixpagename{PHỤ LỤC}
\renewcommand\appendixtocname{PHỤ LỤC}


\addbibresource{Danh_sach_tai_lieu_tham_khao.bib} % chèn file chứa danh mục tài liệu tham khảo vào 

\include{lstlisting} % Phần này cho phép chèn code và formatting code như C, C++, Python

%\makeglossaries
\makenoidxglossaries

% Danh mục thuật ngữ và từ viết tắt — Brightify
\newglossaryentry{VA}{type=\acronymtype, name={V-A},
  description={Valence--Arousal --- hai trục (mức tích cực và mức năng lượng) của mô hình cảm xúc vòng tròn của Russell}}

\newglossaryentry{MER}{type=\acronymtype, name={MER},
  description={Nhận dạng cảm xúc trong âm nhạc (Music Emotion Recognition)}}

\newglossaryentry{SSL}{type=\acronymtype, name={SSL},
  description={Học tự giám sát (Self-Supervised Learning)}}

\newglossaryentry{MuQ}{type=\acronymtype, name={MuQ},
  description={Mô hình biểu diễn âm nhạc tự giám sát dùng Mel Residual Vector Quantization (Zhu và cộng sự, 2025)}}

\newglossaryentry{MERT}{type=\acronymtype, name={MERT},
  description={Mô hình hiểu âm nhạc tự giám sát (Music undERstanding model with large-scale Training)}}

\newglossaryentry{NRCVAD}{type=\acronymtype, name={NRC-VAD},
  description={Từ điển Valence--Arousal--Dominance của NRC (Mohammad, 2018)}}

\newglossaryentry{EWE}{type=\acronymtype, name={EWE},
  description={Bộ ước lượng có trọng số theo độ tin cậy (Evaluator Weighted Estimator)}}

\newglossaryentry{RBF}{type=\acronymtype, name={RBF},
  description={Hàm cơ sở bán kính (Radial Basis Function)}}

\newglossaryentry{MMR}{type=\acronymtype, name={MMR},
  description={Độ liên quan biên cực đại (Maximal Marginal Relevance)}}

\newglossaryentry{CDF}{type=\acronymtype, name={CDF},
  description={Hàm phân phối tích lũy (Cumulative Distribution Function)}}

\newglossaryentry{ICEAS}{type=\acronymtype, name={ICEAS},
  description={Khảo sát liên kết màu sắc--cảm xúc quốc tế (Jonauskaite và cộng sự, 2020)}}

\newglossaryentry{NDCG}{type=\acronymtype, name={NDCG},
  description={Độ lợi tích lũy chiết khấu chuẩn hóa (Normalized Discounted Cumulative Gain)}}

\newglossaryentry{FDR}{type=\acronymtype, name={FDR},
  description={Tỉ lệ phát hiện sai (False Discovery Rate)}}

\newglossaryentry{TE}{type=\acronymtype, name={TE},
  description={Sai số nhắm trúng đích trong không gian V-A (Targeting Error)}}

\newglossaryentry{API}{type=\acronymtype, name={API},
  description={Giao diện lập trình ứng dụng (Application Programming Interface)}}

\newglossaryentry{SPA}{type=\acronymtype, name={SPA},
  description={Ứng dụng trang đơn (Single-Page Application)}}

\newglossaryentry{LLM}{type=\acronymtype, name={LLM},
  description={Mô hình ngôn ngữ lớn (Large Language Model)}}

\newglossaryentry{BPM}{type=\acronymtype, name={BPM},
  description={Số nhịp trên phút --- thước đo nhịp độ bản nhạc (Beats Per Minute)}}

\newglossaryentry{Oklab}{type=\acronymtype, name={Oklab},
  description={Không gian màu đồng đều theo tri giác (Ottosson, 2020)}}

\newglossaryentry{CLAP}{type=\acronymtype, name={CLAP},
  description={Mô hình tương phản ngôn ngữ--âm thanh (Contrastive Language--Audio Pretraining)}}

\newglossaryentry{LUFS}{type=\acronymtype, name={LUFS},
  description={Đơn vị đo độ to chuẩn hóa (Loudness Units relative to Full Scale)}}

\newglossaryentry{ICC}{type=\acronymtype, name={ICC},
  description={Hệ số tương quan nội lớp, đo độ tin cậy test--retest (Intraclass Correlation Coefficient)}}

\newglossaryentry{KS}{type=\acronymtype, name={KS},
  description={Kiểm định Kolmogorov--Smirnov so sánh hai phân bố}}

\newglossaryentry{LOOCV}{type=\acronymtype, name={LOO-CV},
  description={Kiểm chứng chéo bỏ-một-ra (Leave-One-Out Cross-Validation)}}

\newglossaryentry{DEAM}{type=\acronymtype, name={DEAM},
  description={Tập dữ liệu cảm xúc âm nhạc có nhãn người, dùng cho hồi quy Arousal (Database for Emotional Analysis of Music)}}

\newglossaryentry{PMEmo}{type=\acronymtype, name={PMEmo},
  description={Tập dữ liệu cảm xúc âm nhạc dùng để kiểm chuyển miền (Popular Music Emotion dataset)}}

\newglossaryentry{XLMR}{type=\acronymtype, name={XLM-R},
  description={Mô hình ngôn ngữ đa ngữ RoBERTa (XLM-RoBERTa, Conneau và cộng sự, 2020)}}

\newglossaryentry{ViSoBERT}{type=\acronymtype, name={ViSoBERT},
  description={Mô hình ngôn ngữ tiền huấn luyện cho mạng xã hội tiếng Việt (Nguyễn và cộng sự, 2023)}}

\newglossaryentry{REST}{type=\acronymtype, name={REST},
  description={Phong cách kiến trúc giao diện lập trình ứng dụng trên HTTP (Representational State Transfer)}}

\newglossaryentry{JSON}{type=\acronymtype, name={JSON},
  description={Định dạng trao đổi dữ liệu dạng văn bản (JavaScript Object Notation)}}


% ===================================================


\fancypagestyle{plain}{%
\fancyhf{} % clear all header and footer fields
\fancyfoot[RO,RE]{\thepage} %RO=right odd, RE=right even
\renewcommand{\headrulewidth}{0pt}
\renewcommand{\footrulewidth}{0pt}}

\setlength{\headheight}{10pt}

\def \TITLE{ĐỒ ÁN TỐT NGHIỆP}
\def \AUTHOR{NGUYỄN VĂN ABC}

% ===================================================
\titleformat{\chapter}[hang]{\centering\bfseries}{CHƯƠNG \thechapter.\ }{0pt}{}[]

\titleformat 
    {\chapter} % command
    [hang] % shape
    {\centering\bfseries} % format
    {CHƯƠNG \thechapter.\ } % label
    {0pt} %sep
    {} % before
    [] % after
\titlespacing*{\chapter}{0pt}{-20pt}{20pt}

\titleformat
    {\section} % command
    [hang] % shape
    {\bfseries} % format
    {\thechapter.\arabic{section}\ \ \ \ } % label
    {0pt} %sep
    {} % before
    [] % after
\titlespacing{\section}{0pt}{\parskip}{0.5\parskip}

\titleformat
    {\subsection} % command
    [hang] % shape
    {\bfseries} % format
    {\thechapter.\arabic{section}.\arabic{subsection}\ \ \ \ } % label
    {0pt} %sep
    {} % before
    [] % after
\titlespacing{\subsection}{30pt}{\parskip}{0.5\parskip}

\renewcommand\thesubsubsection{\alph{subsubsection}}
\titleformat
    {\subsubsection} % command
    [hang] % shape
    {\bfseries} % format
    {\alph{subsubsection}, \ } % label
    {0pt} %sep
    {} % before
    [] % after
\titlespacing{\subsubsection}{50pt}{\parskip}{0.5\parskip}

% \newcommand{\titlesize}{\fontsize{18pt}{23pt}\selectfont}
% \newcommand{\subtitlesize}{\fontsize{16pt}{21pt}\selectfont}
% \titleclass{\part}{top}
% \titleformat{\part}[display]
%   {\normalfont\huge\bfseries}{\centering}{20pt}{\Huge\centering}
% \titlespacing{\part}{0pt}{em}{1em}
% \titlespacing{\section}{0pt}{\parskip}{0.5\parskip}
% \titlespacing{\subsection}{0pt}{\parskip}{0.5\parskip}
% \titlespacing{\subsubsection}{0pt}{\parskip}{0.5\parskip}



% ===================================================
\usepackage{hyperref}
\hypersetup{pdfborder = {0 0 0}} %
\hypersetup{pdftitle={\TITLE},
	pdfauthor={\AUTHOR}}
	
\usepackage[all]{hypcap} % Cho phép tham chiếu chính xác đến hình ảnh và bảng biểu

\graphicspath{{figures/}{../figures/}} % Thư mục chứa các hình ảnh

\counterwithin{figure}{chapter} % Đánh số hình ảnh kèm theo chapter. Ví dụ: Hình 1.1, 1.2,..

\title{\bf \TITLE}
\author{\AUTHOR}

\setcounter{secnumdepth}{3} % Cho phép subsubsection trong report
% \setcounter{tocdepth}{3} % Chèn subsubsection vào bảng mục lục

\theoremstyle{definition}
\newtheorem{example}{Ví dụ}[chapter] % Định nghĩa môi trường ví dụ

\onehalfspacing
% Tham số float: cho phép hình chiếm nhiều phần trang, ít text tối thiểu -> lấp đầy trang, tránh trang float gần trống và khoảng trắng lớn
\renewcommand{\topfraction}{0.9}
\renewcommand{\bottomfraction}{0.9}
\renewcommand{\textfraction}{0.06}
\renewcommand{\floatpagefraction}{0.8}
\setcounter{topnumber}{3}
\setcounter{bottomnumber}{2}
\setcounter{totalnumber}{4}
%Khoảng cách xuống dòng
\setlength{\parskip}{6pt}
%Lùi đầu dòng
\setlength{\parindent}{15pt}



% =========================== BODY ===============
\begin{document}
% \newgeometry{top=2cm, bottom=2cm, left=2cm, right=2cm}
\subfile{Bia} % Phần bìa
\subfile{Bia_lot} % Phần bìa
% \restoregeometry

% ===================================================
\pagenumbering{roman}
% \pagestyle{empty} % Header và footer rỗng
%\newpage
%\subfile{chapters/0_1_subject.tex}

\newpage
\pagenumbering{gobble}
\subfile{Chuong/0_2_Loi_cam_on.tex}

\newpage
\pagenumbering{gobble}
\subfile{Chuong/0_3_Tom_tat_noi_dung.tex}

% (Tóm tắt tiếng Anh đã tạm bỏ theo yêu cầu)

% ===================================================
% \pagestyle{empty} % Header và footer rỗng
\newpage
\pagenumbering{roman} % Xóa page numbering ở cuối trang
\renewcommand*\contentsname{MỤC LỤC}

\titlecontents{chapter}
    [0.0cm]             % left margin
    {\bfseries\vspace{0.3cm}}                  % above code
    {{\bfseries{\scshape}
    CHƯƠNG \thecontentslabel.\ }}
    % numbered format
    {}         % unnumbered format
    {\titlerule*[0.3pc]{.}\contentspage}         % filler-page-format, e.g dots

    
\titlecontents{section}
    [0.0cm]             % left margin
    {\vspace{0.3cm}}                  % above code
    {\thecontentslabel \ } % numbered format
    {}         % unnumbered format
    {\titlerule*[0.3pc]{.}\contentspage}         % filler-page-format, e.g dots
    
\titlecontents{subsection}
    [1.0cm]             % left margin
    {\vspace{0.3cm}}                  % above code
    {\thecontentslabel \ } % numbered format
    {}         % unnumbered format
    {\titlerule*[0.3pc]{.}\contentspage}         % filler-page-format, e.g dots

 % Tạo mục lục tự động
\addtocontents{toc}{\protect\thispagestyle{empty}}
\tableofcontents 
\thispagestyle{empty}
\cleardoublepage

% \pagenumbering{roman}
%Tạo danh mục hình vẽ.
\renewcommand{\listfigurename}{DANH MỤC HÌNH VẼ}
{\let\oldnumberline\numberline
\renewcommand{\numberline}{Hình~\oldnumberline}
\listoffigures} 
% \phantomsection\addcontentsline{toc}{section}{\numberline {} DANH MỤC HÌNH VẼ}
\newpage


 %Tạo danh mục bảng biểu.
\renewcommand{\listtablename}{DANH MỤC BẢNG BIỂU}
{\let\oldnumberline\numberline
\renewcommand{\numberline}{Bảng~\oldnumberline}
\listoftables}
% \phantomsection\addcontentsline{toc}{section}{\numberline {} DANH MỤC BẢNG BIỂU}

\glsaddall 
 \renewcommand*{\glossaryname}{Danh sách thuật ngữ}
\renewcommand*{\acronymname}{DANH MỤC THUẬT NGỮ VÀ TỪ VIẾT TẮT}
\renewcommand*{\entryname}{Thuật ngữ}
\renewcommand*{\descriptionname}{Ý nghĩa}
\printnoidxglossaries
 %\phantomsection\addcontentsline{toc}{section}{\numberline {} DANH MỤC THUẬT NGỮ VÀ TỪ VIẾT TẮT}

% \newpage
 %\subfile{chapters/0_5_Danh_muc_viet_tat.tex}

% \newpage
% \subfile{chapters/0_6_Thuat_ngu.tex}
% ===================================================


\newpage
\pagenumbering{arabic}

\pagestyle{fancy}
\fancyhf{}
\fancyhead[RE, LO]{\leftmark}
%\fancyhead[LE]{\rightmark}
\fancyfoot[RE, LO]{\thepage}

\chapter{GIỚI THIỆU ĐỀ TÀI}
\label{chapter:Introduction}
\subfile{Chuong/1_Gioi_thieu} % Phần mở đầu

\newpage
%\pagestyle{fancy} % Áp dụng header và footer
\chapter{KHẢO SÁT VÀ PHÂN TÍCH YÊU CẦU}
\label{chapter:Related_works}
\subfile{Chuong/2_Khao_sat}


\newpage
%\pagestyle{fancy} % Áp dụng header và footer
\chapter{NỀN TẢNG LÝ THUYẾT VÀ CÔNG NGHỆ SỬ DỤNG}
\label{chapter:Methodology}
\subfile{Chuong/3_Cong_nghe}

\newpage
%\pagestyle{fancy} % Áp dụng header và footer
\chapter{PHÂN TÍCH VÀ THIẾT KẾ HỆ THỐNG}
\label{chapter:Design}
\subfile{Chuong/4_Thiet_ke}

\newpage
%\pagestyle{fancy} % Áp dụng header và footer
\chapter{XÂY DỰNG VÀ TRIỂN KHAI HỆ THỐNG}
\label{chapter:Implementation}
\subfile{Chuong/5_Xay_dung}

\newpage
%\pagestyle{fancy} % Áp dụng header và footer
\chapter{THỬ NGHIỆM VÀ ĐÁNH GIÁ}
\label{chapter:Evaluation}
\subfile{Chuong/6_Danh_gia}

\newpage
%\pagestyle{fancy} % Áp dụng header và footer
\chapter{CÁC GIẢI PHÁP VÀ ĐÓNG GÓP NỔI BẬT}
\label{chapter:SolutionAndContribution}
\subfile{Chuong/5_Giai_phap_dong_gop}
\newpage
%\pagestyle{fancy} % Áp dụng header và footer
\chapter{KẾT LUẬN VÀ HƯỚNG PHÁT TRIỂN} %Kết luận và hướng phát triển}
\label{chapter:conclusion}
\subfile{Chuong/6_Ket_luan}

% (Mục "Một số lưu ý về tài liệu tham khảo" đã xóa theo yêu cầu)


% ===================================================
\newpage
\renewcommand\bibname{TÀI LIỆU THAM KHẢO}
\printbibliography
\phantomsection\addcontentsline{toc}{chapter}{TÀI LIỆU THAM KHẢO}

\appendixpage
\appendices
\addappheadtotoc

% \chapter*{PHỤ LỤC} %Kết luận và hướng phát triển}

%\mainmatter
\titleformat{\chapter}[hang]{\centering\bfseries}{ \thechapter.\ }{0pt}{}[]
\titlespacing*{\chapter}{0pt}{-20pt}{20pt}

\titlecontents{chapter}
    [0.0cm]             % left margin
    {\bfseries\vspace{0.3cm}}                  % above code
    {{\bfseries{\scshape} \thecontentslabel.\ }} % numbered format
    {}         % unnumbered format
    {\titlerule*[0.3pc]{.}\contentspage}         % filler-page-format, e.g dots
\chapter{ĐẶC TẢ USE CASE}
\subfile{Chuong/Phu_luc_B}
% (Các phụ lục B/C/D — API, tham số, lược đồ kho — tạm bỏ theo yêu cầu, chỉ giữ Phụ lục A)
\end{document}
